\chapter{Pricing and Numerics}\label{ch:iv2:pricing-and-numerics}

\begin{ivtrains}
Deriving a price three ways, pricing what has no formula, and knowing how wrong a tree, a grid or a simulation is, and why. Unless a stem says otherwise, the model is Black-Scholes, with a constant rate and a constant volatility and no dividends.
\end{ivtrains}

\section{Warm-ups}

\begin{question}[$\star$ \iqroles{QR, QT, QD} \iqfirm{any} \iqlevel{junior}]\label{iq:iv2:pricing-and-numerics:1}
Write down the Black-Scholes price of a call. With a stock and strike at 100, 20\% volatility, one year and a zero rate, give the price in your head by a rule of thumb, and say what the exact value is. Do the same at 30\% and three months.
\end{question}

\begin{question}[$\star$ \iqroles{QR, QT, QD} \iqfirm{any} \iqlevel{junior}]\label{iq:iv2:pricing-and-numerics:2}
Two stocks are both at 100, with the same volatility. Everyone expects one to return 20\% a year, and the other 2\%. The rate is the same for both. Which has the dearer one-year call in the Black-Scholes model? Why does a trader feel this must be wrong, and what do you answer?
\end{question}

\begin{question}[$\star$ \iqroles{QR, QT, QD} \iqfirm{any} \iqlevel{junior}]\label{iq:iv2:pricing-and-numerics:3}
A stock at 100 goes to 120 or to 90 tomorrow, and the rate is zero. Price a call struck at 100 by building a portfolio of stock and cash that pays the same in both states. What is the risk-neutral chance of going up? What changes if the true chance of going up is 80\%?
\end{question}

\begin{question}[$\star$ \iqroles{QR, QD, Risk} \iqfirm{any} \iqlevel{junior}]\label{iq:iv2:pricing-and-numerics:4}
Your payoff has a standard deviation of \$15. With 10\,000 Monte Carlo paths, what is the standard error of the price? How many paths do you need for a standard error of 1 cent?
\end{question}

\begin{question}[$\star$ \iqroles{QR, QT, QD} \iqfirm{any} \iqlevel{junior}]\label{iq:iv2:pricing-and-numerics:5}
A one-year call is worth \$8. The same call that dies if the stock ever trades below 90 is worth \$6. What is the same call worth if it only comes alive once the stock trades below 90? Why is this true in any model? What does the \$6 tend to if the barrier falls toward zero?
\end{question}

\begin{question}[$\star$ \iqroles{QR, QT, QD} \iqfirm{any} \iqlevel{junior}]\label{iq:iv2:pricing-and-numerics:6}
A contract pays \$1 if a stock finishes above 100 in a year, and it is worth \$0.46. The rate is 5\%, compounded continuously. What is the contract that pays \$1 if the stock finishes below 100 worth? What assumption do you need about the stock? Which volatility did you use?
\end{question}

\begin{question}[$\star$ \iqroles{QR, QT, QD} \iqfirm{any} \iqlevel{junior}]\label{iq:iv2:pricing-and-numerics:7}
An at-the-money one-year option pays on the continuous average of a stock over the year, and not on its final price. The stock is at 100 with 30\% volatility and a zero rate. Is it dearer or cheaper than the ordinary call? Roughly what volatility should you use for the average, and what is the price by the rule of thumb?
\end{question}

\begin{question}[$\star$ \iqroles{QD, QR} \iqfirm{any} \iqlevel{junior}]\label{iq:iv2:pricing-and-numerics:8}
You build a binomial tree with 20 time steps. How many nodes does it have if up then down lands on the same price as down then up? How many if it does not? Why does this matter for barrier-free, path-independent options, and which options lose this saving?
\end{question}

\begin{question}[$\star$ \iqroles{QR, QD} \iqfirm{any} \iqlevel{junior}]\label{iq:iv2:pricing-and-numerics:9}
You use antithetic paths: every normal draw $Z$ is used with $-Z$. What happens when you estimate $E[Z]$? What happens for $E[Z^2]$? For which kinds of payoff does the trick help?
\end{question}

\begin{question}[$\star$ \iqroles{QR, QD} \iqfirm{any} \iqlevel{junior}]\label{iq:iv2:pricing-and-numerics:10}
You have two independent standard normals, $Z_1$ and $Z_2$. Build two normals with a correlation of 0.6. What do you get for the second if $Z_1 = 1$ and $Z_2 = -1$? What does this have to do with a Cholesky factor?
\end{question}

\begin{question}[$\star$ \iqroles{QR, QT, QD, Risk} \iqfirm{any} \iqlevel{junior}]\label{iq:iv2:pricing-and-numerics:11}
A stock is at 100, pays no dividends, and the rate is 5\%, compounded continuously. What is the fair one-year forward price? You hold a contract to buy at 100 in one year. What is it worth today? Why do you need no volatility?
\end{question}

\begin{question}[$\star$ \iqroles{QD, QR} \iqfirm{any} \iqlevel{junior}]\label{iq:iv2:pricing-and-numerics:12}
You solve the pricing equation with an explicit finite-difference grid, in time and one price. You halve the price step. By what factor does the work go up, and why not just by two?
\end{question}

\begin{question}[$\star$ \iqroles{QD, QR} \iqfirm{any} \iqlevel{junior}]\label{iq:iv2:pricing-and-numerics:13}
You can draw uniform random numbers. How do you draw an exponential with rate 2? What value do you get from $U = 0.5$? What is the general method, and why is it not the usual way of drawing a normal?
\end{question}

\section{Deriving prices}

\begin{question}[$\star\star$ \iqroles{QR, QT, QD} \iqfirm{any} \iqlevel{junior}]\label{iq:iv2:pricing-and-numerics:14}
A stock is at 110, with 30\% volatility, half a year to expiry and a zero rate. Under the risk-neutral measure, what is the chance it finishes above 100? Derive it from the distribution of the log price, and do not just quote $\Phi(d_2)$.
\end{question}

\begin{question}[$\star\star$ \iqroles{QR, QD} \iqfirm{any} \iqlevel{junior, experienced}]\label{iq:iv2:pricing-and-numerics:15}
Derive the Black-Scholes call price as a risk-neutral expectation, with the rate at zero. In particular, how do you handle $E[S_T \mathbf{1}\{S_T > K\}]$? Check it for a \$110 stock, a 100 strike, 30\% volatility and half a year.
\end{question}

\begin{question}[$\star\star$ \iqroles{QR, QD} \iqfirm{any} \iqlevel{junior, experienced}]\label{iq:iv2:pricing-and-numerics:16}
Derive the Black-Scholes equation with a delta-hedged portfolio. Where has the stock's expected return gone? For an at-the-money one-year call on a \$100 stock, with 20\% volatility and a zero rate, the gamma is 0.0198. What must the theta be?
\end{question}

\begin{question}[$\star\star$ \iqroles{QR, QD} \iqfirm{hedge fund, bank, market maker} \iqlevel{junior, experienced}]\label{iq:iv2:pricing-and-numerics:17}
A contract pays $S_T^2$ in one year, on a stock at \$100 with 20\% volatility and a zero rate. Solve for its value by trying a solution of the form $V = S^2 g(t)$ in the Black-Scholes equation. What is $g$, and what is the value? Check it with an expectation.
\end{question}

\begin{question}[$\star\star$ \iqroles{QR, QT, QD} \iqfirm{any} \iqlevel{junior}]\label{iq:iv2:pricing-and-numerics:18}
Show that a call is made of an asset-or-nothing option, which pays the stock if it finishes above the strike, and a cash-or-nothing option. What is each worth for a stock and strike of 100, 20\% volatility, a zero rate and one year? Why is the asset one's value not $S$ times the cash one's probability?
\end{question}

\begin{question}[$\star\star$ \iqroles{QR, QT, QD} \iqfirm{bank, hedge fund, market maker} \iqlevel{junior, experienced}]\label{iq:iv2:pricing-and-numerics:19}
With a normal (Bachelier) model in place of a lognormal one, a stock at 100 has a dollar volatility of \$20 a year. With a zero rate, what is the at-the-money one-year call worth? Compare it with a lognormal model at 20\%. What implied lognormal volatilities does the normal model give at strikes of 80 and 120? Why does it have a skew?
\end{question}

\begin{question}[$\star\star$ \iqroles{QR, QT} \iqfirm{any} \iqlevel{junior, experienced}]\label{iq:iv2:pricing-and-numerics:20}
A stock with no drift moves as a Brownian motion, with a dollar volatility of \$20 a year, starting at 100. What is the chance it touches 120 in a year? Use the reflection argument. If you check only closing prices, on 252 days, what is the chance, roughly?
\end{question}

\begin{question}[$\star\star$ \iqroles{QR, QT, QD} \iqfirm{bank, hedge fund, market maker} \iqlevel{junior, experienced}]\label{iq:iv2:pricing-and-numerics:21}
A contract pays \$1 the first time a stock reaches 110, and it never expires. The stock is at 100, a lognormal martingale with 20\% volatility, and the rate is zero. What is it worth? Now it lasts a year: a desk's rule of thumb is twice the European digital at 110. What is that, and what is the exact answer?
\end{question}

\begin{question}[$\star\star$ \iqroles{QR, QT} \iqfirm{hedge fund, market maker, bank} \iqlevel{junior, experienced}]\label{iq:iv2:pricing-and-numerics:22}
A contract pays $\ln(S_T/S_0)$ at expiry. The stock has 20\% volatility, a zero rate and one year to go. What is it worth? Why is it negative, when the stock is a fair bet? What do you get if you sell the contract and hold $1/S$ shares, rebalancing as you go?
\end{question}

\begin{question}[$\star\star$ \iqroles{QR, QT, QD} \iqfirm{bank, hedge fund, market maker} \iqlevel{junior, experienced}]\label{iq:iv2:pricing-and-numerics:23}
In six months, you will be given a six-month at-the-money call with a strike set at that day's stock price. The stock is at 100, with 20\% volatility and a zero rate. What is it worth today? Why does the price not depend on the stock's level in six months?
\end{question}

\begin{question}[$\star\star$ \iqroles{QR, QT, QD} \iqfirm{bank, hedge fund, market maker} \iqlevel{junior, experienced}]\label{iq:iv2:pricing-and-numerics:24}
You hold a contract that lets you decide in six months whether it is a call or a put, both struck at 100 and both expiring in a year. The stock is at 100, with 20\% volatility and a zero rate. Show that it is worth the same as a one-year call plus a six-month put. What is the value? How does that compare with a one-year straddle?
\end{question}

\section{Discretising}

\begin{question}[$\star\star$ \iqroles{QD, QR} \iqfirm{any} \iqlevel{junior, experienced}]\label{iq:iv2:pricing-and-numerics:25}
You simulate a lognormal stock by Euler, $S_{i+1} = S_i(1 + \sigma\sqrt{\Delta t}\,Z)$, with $\sigma = 100\%$ over one year and four steps, and no drift. What is the chance of a negative price at one step? What is $E[S_T^2]/S_0^2$ in this scheme, and what is the true value? What should you simulate instead?
\end{question}

\begin{question}[$\star\star$ \iqroles{QD, QR} \iqfirm{any} \iqlevel{junior, experienced}]\label{iq:iv2:pricing-and-numerics:26}
Give the standard binomial tree's parameters for 20\% volatility, with a zero rate. For $\Delta t = 0.25$, what are $u$, $d$ and the probability of an up-move? For an at-the-money one-year call, the exact value is 7.9656. The tree gives 7.7692 with 10 steps, 8.1472 with 11, 7.9457 with 100, and 7.9852 with 101. Why does it oscillate? What cheap fix do these numbers suggest?
\end{question}

\begin{question}[$\star\star$ \iqroles{QD, QR} \iqfirm{any} \iqlevel{junior, experienced}]\label{iq:iv2:pricing-and-numerics:27}
You use an explicit finite-difference scheme on prices from 0 to 200, in steps of \$1, for a stock with 20\% volatility and a zero rate. What is the largest stable time step, and how many steps for a year? What happens if you go a little beyond it? Why is this a reason to go implicit?
\end{question}

\begin{question}[$\star\star$ \iqroles{QR, QD, Risk} \iqfirm{any} \iqlevel{junior, experienced}]\label{iq:iv2:pricing-and-numerics:28}
You are given a correlation of 0.9 between A and B, and 0.9 between A and C, and someone sets the correlation of B and C to zero. What goes wrong when you take a Cholesky factor? Show it with a number. What is the lowest correlation of B and C that is allowed? How would you repair a matrix like that in practice?
\end{question}

\begin{question}[$\star\star$ \iqroles{QD, QR, QT} \iqfirm{any} \iqlevel{junior, experienced}]\label{iq:iv2:pricing-and-numerics:29}
An at-the-money one-year call on a \$100 stock with a zero rate trades at \$12. Find the implied volatility using a first guess from the rule of thumb, then Newton's method. How many steps do you need for six digits? Where does Newton's method fail in practice, and what is the safeguard?
\end{question}

\begin{question}[$\star\star$ \iqroles{QD, QR, Risk} \iqfirm{any} \iqlevel{junior, experienced}]\label{iq:iv2:pricing-and-numerics:30}
You estimate an at-the-money call's delta on a \$100 stock, with 20\% volatility, one year and a zero rate, by bumping the stock \$1 up and \$1 down. You use 100\,000 paths. The payoff's standard deviation is about \$13.15. What is the standard error of the delta if the up-run and the down-run use different random numbers? What about the same random numbers? Roughly what does that do?
\end{question}

\begin{question}[$\star\star$ \iqroles{QR, QT, QD} \iqfirm{any} \iqlevel{junior, experienced}]\label{iq:iv2:pricing-and-numerics:31}
A basket holds equal weights in two stocks at 100, both with 30\% volatility and with a correlation of 0.5. What is the basket's volatility? Price an at-the-money one-year call on it, with a zero rate, using a single lognormal. How would you check that? What is the call worth with a correlation of 0 and of 1?
\end{question}

\begin{question}[$\star\star$ \iqroles{QR, QD} \iqfirm{any} \iqlevel{junior, experienced}]\label{iq:iv2:pricing-and-numerics:32}
You price an at-the-money call by simulation, and you use the final stock price as a control variate, since you know its mean. If the correlation between the call's payoff and the stock is 0.89, by how much does the variance fall? What is the optimal multiplier, and how do you estimate it? How would the correlation change for a deep in-the-money call?
\end{question}

\section{Exotics and schemes}

\begin{question}[$\star\star\star$ \iqroles{QR, QD} \iqfirm{any} \iqlevel{junior, experienced}]\label{iq:iv2:pricing-and-numerics:33}
An option pays on the continuous geometric average of a stock over one year. The stock and strike are at 100, with 30\% volatility and a zero rate. Show that the average is lognormal, and find its volatility and forward. What is the price? Why is there a closed form here, and not for an arithmetic average?
\end{question}

\begin{question}[$\star\star\star$ \iqroles{QR, QD} \iqfirm{any} \iqlevel{junior, experienced}]\label{iq:iv2:pricing-and-numerics:34}
An option pays on the continuous arithmetic average of a stock over one year. The stock and strike are at 100, with 30\% volatility and a zero rate. Match a lognormal to the first two moments. The second moment of the continuous average is $E[A^2] = 2S^2(e^{\sigma^2T} - 1 - \sigma^2T)/(\sigma^2T)^2$. What is the price? How do you check it by simulation with few paths, and what does an option on the geometric average do for you? By what factor does the standard error fall if the two payoffs correlate at 0.9992?
\end{question}

\begin{question}[$\star\star\star$ \iqroles{QR, QT, QD} \iqfirm{any} \iqlevel{junior, experienced}]\label{iq:iv2:pricing-and-numerics:35}
An option lets you swap stock B for stock A at expiry, paying $(S_A - S_B)^+$. Both are at 100. Their volatilities are 30\% and 20\%, the correlation is 0.5, and it is a one-year option. What is it worth? Show why the rate does not enter. What is the trick? Check it in the limit where B has no volatility.
\end{question}

\begin{question}[$\star\star\star$ \iqroles{QR, QT} \iqfirm{bank, hedge fund, market maker} \iqlevel{junior, experienced}]\label{iq:iv2:pricing-and-numerics:36}
A one-year call struck at 100 dies if the stock ever trades at 100 or lower. The stock is at 110, with a zero rate. What is it worth, in a model with any continuous martingale for the stock? Why does volatility not come in? Compare it with an ordinary call at 30\% volatility, and say what the difference is worth to someone.
\end{question}

\begin{question}[$\star\star\star$ \iqroles{QR, QT, QD} \iqfirm{bank, hedge fund, market maker} \iqlevel{junior, experienced}]\label{iq:iv2:pricing-and-numerics:37}
A one-year call struck at 100 dies if the stock ever trades at 130. The stock is at 100, with a zero rate. Use the method of images to price it, taking 20\%, 15\% and 25\% for the volatility. What is the sign of its vega? How does this compare with an ordinary call, and why?
\end{question}

\begin{question}[$\star\star\star$ \iqroles{QR, QD, QT} \iqfirm{bank, hedge fund} \iqlevel{junior, experienced}]\label{iq:iv2:pricing-and-numerics:38}
A one-year up-and-out call, on a stock at 100, is struck at 100 with a barrier at 130, 20\% volatility and a zero rate. The barrier is only watched at the close, on 252 days. Its continuous price is 2.97. Should the daily one be higher or lower? What is the standard shift to the barrier that fixes the formula, and what does it give? How large is the effect on the barrier level?
\end{question}

\begin{question}[$\star\star\star$ \iqroles{QR, QD} \iqfirm{bank, hedge fund} \iqlevel{junior, experienced}]\label{iq:iv2:pricing-and-numerics:39}
You price a one-year up-and-out call with a continuously watched barrier at 130, with the stock and strike at 100, 20\% volatility and a zero rate. With monthly Monte Carlo steps, you get 3.73. The right answer is 2.97. What is the bias? How do you fix it without more steps? Given a month that starts at 120 and ends at 125, with 20\% volatility, what is the chance it touched 130 in between?
\end{question}

\begin{question}[$\star\star\star$ \iqroles{QR, QD, Risk} \iqfirm{any} \iqlevel{experienced}]\label{iq:iv2:pricing-and-numerics:40}
You want a call's delta by simulation, for a \$100 stock and strike, 20\% volatility, one year and a zero rate. Derive the pathwise estimate and the likelihood-ratio estimate. What are their means and standard deviations per path? Which do you use for a digital, and why?
\end{question}

\begin{question}[$\star\star\star$ \iqroles{QR, QT} \iqfirm{any} \iqlevel{junior, experienced}]\label{iq:iv2:pricing-and-numerics:41}
A stock follows a Brownian motion with no drift and a dollar volatility of \$20 a year. A contract pays the difference between its highest price in the year and its starting price, and it is watched continuously. What is it worth, by reflection? How does that compare with an at-the-money straddle in the same model? What would you expect with daily checks only?
\end{question}

\begin{question}[$\star\star\star$ \iqroles{QR, QT, QD} \iqfirm{market maker, bank, hedge fund} \iqlevel{junior, experienced}]\label{iq:iv2:pricing-and-numerics:42}
A \$1 digital pays above 100 in a year. The stock is at 100 and the rate is zero. At-the-money implied volatility is 20\%, the vega is 39.7 per unit of volatility, and the smile falls 0.2 points of volatility for each \$1 of strike. What is the digital worth with no skew, and with it? Derive the correction. What is the sign for most equity indexes? How much does it change the chance?
\end{question}

\begin{question}[$\star\star\star$ \iqroles{QR, QD} \iqfirm{any} \iqlevel{experienced}]\label{iq:iv2:pricing-and-numerics:43}
You want the chance that a standard normal exceeds 4. With a million plain Monte Carlo draws, what is the relative standard error? Shift the draws by 4 and reweight. What is the weight, and what is the relative standard error now? What can go wrong with a shift when you do it in many dimensions?
\end{question}

\begin{question}[$\star\star\star$ \iqroles{QR, QD} \iqfirm{any} \iqlevel{experienced}]\label{iq:iv2:pricing-and-numerics:44}
You price a one-year at-the-money call by one normal draw per path, $2^{14}$ paths, with 20\% volatility. Compare the error of plain Monte Carlo with a scrambled Sobol sequence, as measured over twenty seeds. Why do the Sobol points do better? In what situations do you lose most of this advantage?
\end{question}

\begin{question}[$\star\star\star$ \iqroles{QR, QD} \iqfirm{any} \iqlevel{junior, experienced}]\label{iq:iv2:pricing-and-numerics:45}
You price a one-year American put on a stock and strike at 100, with 20\% volatility and a 5\% rate. A very large tree gives 6.0903, and the European is 5.5735. A 100-step tree gives 6.0824, a 101-step tree 6.1046 and a 200-step tree 6.0864. The 100-step European tree gives 5.5536. Compare the plain value, the average of 100 and 101 steps, the European control and doubling with extrapolation. Which is best here? Which is the usual textbook choice? What is the stock price below which you exercise today?
\end{question}

\begin{question}[$\star\star\star$ \iqroles{QD, QR} \iqfirm{any} \iqlevel{experienced}]\label{iq:iv2:pricing-and-numerics:46}
You solve a call with five days to go, $T = 0.02$, by Crank-Nicolson on a fine price grid in steps of \$0.25, with five time steps. The price is close enough, but the gamma at the strike comes out as $-1.36$. The true gamma is 0.141. What is happening, even though the scheme is unconditionally stable? What is the standard fix? What does it do to this gamma?
\end{question}

\begin{question}[$\star\star\star$ \iqroles{QR, QD, Risk} \iqfirm{any} \iqlevel{junior, experienced}]\label{iq:iv2:pricing-and-numerics:47}
Find the payoff's exact standard deviation for a one-year call on a \$100 stock with 20\% volatility and a zero rate, with $E[((S_T - K)^+)^2]$ in closed form, for strikes of 100 and 130. How many paths do you need for a 1-cent standard error? How does the error compare with the price at each strike? What does this say about pricing tail options?
\end{question}

\section{Three to build}

\begin{question}[$\star\star\star\star$ \iqroles{QR, QD} \iqfirm{bank, hedge fund, market maker} \iqlevel{experienced}]\label{iq:iv2:pricing-and-numerics:48}
Describe, step by step, the Longstaff-Schwartz regression method for a Bermudan put with 12 monthly exercise dates, struck at 100 on a \$100 stock, with 20\% volatility and a 5\% rate. A tree gives 6.04. What does the method give with 100\,000 paths, and what with 400? Which way is it biased when you reuse the fitting paths, and which way on fresh paths? Why do you regress only on in-the-money paths?
\end{question}

\begin{question}[$\star\star\star\star$ \iqroles{QD, QR} \iqfirm{bank, hedge fund, market maker} \iqlevel{experienced}]\label{iq:iv2:pricing-and-numerics:49}
Set up an American put as a finite-difference problem on a grid in the stock price, for a \$100 stock struck at 100, with 20\% volatility, a 5\% rate and a year to expiry. Write the condition that holds at every point, and how to solve it at each step. What do you get with a \$1 grid and 100 time steps, against 6.0903? Where is the exercise boundary today? What goes wrong if you solve the linear system and then cut it off at the payoff?
\end{question}

\begin{question}[$\star\star\star\star$ \iqroles{QR, QD} \iqfirm{hedge fund, bank, market maker} \iqlevel{experienced}]\label{iq:iv2:pricing-and-numerics:50}
Given a model's characteristic function for the log price, $\phi(u) = E[e^{iu\ln S_T}]$, derive the call price from two probabilities, and write them as integrals. Test it with Black-Scholes at strikes of 100 and 120, with a \$100 stock, 20\% volatility, a zero rate and one year. What does it give, and how much do you lose to the upper limit of integration? Why would anyone do this if a formula exists?
\end{question}
